% L43_P1_obs_teacher.tex
% Algebra 2  Lesson 4-3, Period 1 (Observation Day)  --  TEACHER PACKET
%
% Visual + structural reference: L35_P3_obs_teacher.docx (admin's 3-column
% Lesson Plan Template).
%
% Two tables:
%   1. Header table: 8 rows x 2 cols (Date/Subject, CCSS, Topic Goals, EQ,
%      Materials, Content Obj + criteria, Language Obj + criteria, IEP/ELL).
%   2. Body table: 5 rows x 4 cols (Lesson Part | Teachers will... |
%      Students will... | Questions to consider...).
%
% Pedagogical structure: 3-Reads routine (T reads -> S read+highlight) on the
% first task (Apply, p2), gradual release on the second task (Performance
% Task, p3). Groups of 2-3 labeled A, B, C. Teacher circulates with iPad
% to identify student work to spotlight; doesn't give answers, surfaces
% group thinking.

\documentclass[10pt]{article}
\usepackage[letterpaper, margin=0.55in, landscape]{geometry}
\usepackage{enumitem}
\usepackage{tabularx}
\usepackage{xltabular}
\usepackage{array}
\usepackage{amsmath}
\usepackage{amssymb}
\usepackage{ragged2e}

\setlength{\parindent}{0pt}
\setlength{\parskip}{2pt}
\renewcommand{\arraystretch}{1.15}

\newcolumntype{Y}{>{\RaggedRight\arraybackslash}X}

\begin{document}

% =====================================================================
% TITLE
% =====================================================================
\begin{center}
{\large\bfseries Lesson 4-3, Period 1 (Observation Day) -- Teacher Plan}     \newline
{\small Multiply and Divide Rational Expressions \quad $\bullet$ \quad
Friday 2026-05-15 \quad $\bullet$ \quad Period F (65 min)}
\end{center}

\vspace{0.4em}

% =====================================================================
% HEADER TABLE  (mirrors admin DOCX header: 8 logical rows, 2 cols)
% =====================================================================
\begin{tabularx}{\linewidth}{|>{\bfseries\RaggedRight\arraybackslash}p{1.7in}|Y|>{\bfseries\RaggedRight\arraybackslash}p{1.7in}|Y|}
  \hline
  Date & 2026-05-15 (Fri)
  & Grade Level & 11--12 \\ \hline
  Subject & Algebra 2 / Period F
  & Section & Topic 4 (Rational Functions) -- Lesson 4-3, Period 1 \\ \hline
\end{tabularx}

\vspace{-1.2pt}
\noindent
\begin{tabularx}{\linewidth}{|>{\bfseries\RaggedRight\arraybackslash}p{1.7in}|Y|}
  \hline
  Common Core State Standards &
    HSA-APR.D.6 (rewrite simple rational expressions in different forms);
    HSA-APR.D.7+ (understand that rational expressions form a system
    closed under operations analogous to rational numbers);
    HSA-SSE.A.2 (use the structure of an expression to identify ways to
    rewrite it). \\ \hline
  Topic Goals &
    Use the structure of rational expressions to rewrite them in
    different forms while tracking domain restrictions.
    Build and interpret rational expressions that arise from real-world
    ratios (geometry, probability, efficiency).
    Recognize that simplifying a rational expression does not change
    its domain restrictions. \\ \hline
  Essential Question &
    What stays the same and what changes when we rewrite a rational
    expression in simplified form? \\ \hline
  Materials &
    Student packet (4 pp.); highlighters (one per student); pencils;
    teacher iPad (for circulating + capturing work); whiteboard +
    projector; visible timer; reference: Savvas SE/TE for 4-3. \\ \hline
\end{tabularx}

\vspace{-1.2pt}
\noindent
\begin{tabularx}{\linewidth}{|>{\bfseries\RaggedRight\arraybackslash}p{1.7in}|Y|>{\bfseries\RaggedRight\arraybackslash}p{1.7in}|Y|}
  \hline
  Content Objective(s) &
    I can identify the values that must be excluded from the domain of a
    rational expression. I can simplify rational expressions by factoring
    and canceling common factors. I can build and interpret rational
    expressions that arise from real-world ratios.
  & Criteria for Success / Formative Assessment &
    On Lupe Carnival Part B, student writes the simplified expression
    $\dfrac{x}{6(x+1)}$ AND states $x \neq 0, -1$ with a physical
    interpretation. On Part C, student substitutes $x = 4$ correctly
    AND makes a defensible fairness judgment.
    \newline\smallskip
    \textbf{Self-Reflection:} ``I can explain why the original denominator
    (not the simplified one) determines the domain.'' \\ \hline
  Language Objective &
    Students will explain why a domain restriction must come from the
    \emph{original} denominator (not the simplified one), verbally and
    in writing, with the support of structured talk-time, models, and
    real-time scaffolding/questioning.
  & Criteria for Success / Formative Assessment &
    Use the printed sentence frames to articulate (a) the simplified
    expression and (b) the excluded values with a reason from the
    original expression. Use the comparison frame ``rational expressions
    are similar to rational numbers because \rule{0.6in}{0.4pt}.''
    \newline\smallskip
    \textbf{Self-Reflection:} ``I used at least one sentence frame today
    to explain my thinking.'' \\ \hline
  Modifications / Support for IEP students &
    Tasks pre-chunked (Part A / B / C scaffolds on packet). Sentence
    frames printed inline. Extended time embedded (groups of 2--3 work
    at own pace inside time block). Highlighters provided for the
    3-Reads. Teacher checks in 1:1 during circulation lap 2.
  & Supports for ELL students &
    3-Reads routine (teacher reads first, then students read +
    highlight). Sentence frames on every long-explanation prompt.
    Visual TikZ diagrams for both modeling tasks. Word bank in
    Share/Summary: \emph{domain restriction, simplified, equivalent,
    closed, undefined}. Group A/B/C structure builds peer support. \\ \hline
\end{tabularx}

\vspace{0.6em}

% =====================================================================
% BODY TABLE  (5 rows x 4 cols: Lesson Part | Teachers | Students | Questions)
% =====================================================================
\begin{xltabular}{\linewidth}{|>{\bfseries\RaggedRight\arraybackslash}p{1.4in}|Y|>{\RaggedRight\arraybackslash}p{1.9in}|>{\itshape\RaggedRight\arraybackslash}p{1.55in}|}
  \hline
  \textbf{Lesson Part} & \textbf{Teachers will\ldots} & \textbf{Students will\ldots} & \textbf{Questions to consider\ldots} \\ \hline
  \endfirsthead

  \hline
  \textbf{Lesson Part} & \textbf{Teachers will\ldots} & \textbf{Students will\ldots} & \textbf{Questions to consider\ldots} \\ \hline
  \endhead

  % ---------- DO NOW ----------
  Do Now (10 min) \par
  \emph{Explore \& Reason -- conceptual gateway, DOK 2}
  &
  Hand out student packet at the door.
  Project the graph of $y = x + 2$ from the packet.
  Read prompt aloud once.
  Set timer for 4 minutes silent independent work on Parts A, B, C.
  Then 4 minutes pair share (turn to elbow partner).
  Cold-call 1--2 pairs to share Part C answer; surface ``what kind of
  function might have a graph with a hole or restriction?''
  \newline\smallskip
  \textbf{Bridge to Apply:} ``When the denominator equals zero, the
  function is undefined. That's a \emph{domain restriction} -- and we'll
  see it everywhere today.''
  &
  Read the prompt silently.     \newline
  Sketch a function whose domain excludes $x = 2$.     \newline
  Write Part C on the printed frame.     \newline
  Share thinking with elbow partner.     \newline
  Volunteer (or be cold-called) to share.
  &
  How can you engage students in the lesson?     \newline
  How can you hook them into rational expressions?     \newline
  How are all adults intentionally planned for? \\ \hline

  % ---------- LAUNCH / APPLY ----------
  Launch / Apply (15 min) \newline
  \emph{\#35 Cylinders -- Make Sense \& Persevere, DOK 2} \newline
  \emph{3-READS ROUTINE (full scaffold)}
  &
  \textbf{3-Reads (2 min total):} \newline
  $\bullet$ \textbf{Read 1 (1 min):} \emph{Teacher reads aloud.} Students put pencils down and listen. Ask: ``What is this problem about?'' \newline
  $\bullet$ \textbf{Read 2 (1 min):} Students re-read silently and \emph{highlight} the formulas, the dimensions of A and B, and the question. Ask: ``What information do we need?''
  \newline\smallskip
  \textbf{Part a (5 min):} students work \emph{individually} on A's V/SA. Teacher circulates with iPad noting interesting work. Gather responses from groups A, B, C in turn. \textbf{Do not give the answer} -- collect thinking, surface differences.
  \newline\smallskip
  \textbf{Part b (5 min):} 2 min individual + 3 min turn-and-talk. Share out + review.
  \newline\smallskip
  \textbf{Part c (3 min):} TPS structure. Surface ratio comparison + design implication using the sentence frame on the page.
  &
  Listen during Read 1 (no writing).     \newline
  Highlight key info during Read 2.     \newline
  Work Part a individually for 5 minutes.     \newline
  Be ready to share when called by group letter (A / B / C).     \newline
  Work Part b individually for 2 minutes, then discuss with elbow partner for 3.     \newline
  Apply same TPS pattern to Part c.     \newline
  Use the sentence frame to share the comparison.
  &
  What will you explicitly model and how?     \newline
  How will you know if students are ready to work alone?     \newline
  How will you elicit student thinking without revealing the answer?     \newline
  How are all adults intentionally planned for? \\ \hline

  % ---------- EXPLORE / PERFORMANCE TASK ----------
  Explore / Performance Task (20 min) \newline
  \emph{\#36 Lupe Carnival -- Model with Mathematics, DOK-3 SPINE} \newline
  \emph{3-READS (gradual release)}
  &
  \textbf{3-Reads -- gradual release (2 min total):} \newline
  $\bullet$ \textbf{Read 1 (1 min):} \emph{Students read silently first} (no teacher read aloud). Ask: ``What's the question?'' \newline
  $\bullet$ \textbf{Read 2 (1 min):} Pair share: ``What information do we need to build the probability?''
  \newline\smallskip
  \textbf{Group work} on Parts A, B, C in groups of 2--3. Teacher circulates with iPad, capturing build strategies, errors, and domain interpretations.
  \newline\smallskip
  \textbf{Mid-explore checkpoint (10 min mark):} brief pause -- ``Is your simplified expression $\dfrac{x}{6(x+1)}$? Anyone get something else?'' Don't reveal whether correct -- surface the discussion.
  \newline\smallskip
  \textbf{End of phase:} cold-call by group letter -- Group A shares Part A reasoning; Group B shares Part B (simplify + domain); Group C shares Part C (fairness judgment). Rotate based on iPad notes.
  &
  Read silently first (Read 1).     \newline
  Discuss with group what the question is asking (Read 2).     \newline
  Build the rational expression as a group on Part A.     \newline
  Simplify and state domain restrictions in physical context (Part B).     \newline
  Substitute $x = 4$ and judge fairness (Part C).     \newline
  Use the printed sentence frames.     \newline
  Share out one part per group letter.
  &
  What will students do for most of the session? What will they practice?     \newline
  What type of academic conversations will students engage in?     \newline
  How will you know which groups need re-direction?     \newline
  How will you press for the DOK-3 lift (build $\rightarrow$ simplify $\rightarrow$ judge) without revealing the path? \\ \hline

  % ---------- SHARE / SUMMARY ----------
  Share / Summary (5 min) \par
  \emph{Synthesis around throughline}
  &
  Synthesize the day's throughline: \par
  \emph{``Domain restriction is the soul of a rational expression --
  it survives even when factors cancel.''}     \newline
  Connect Do Now $\rightarrow$ Apply $\rightarrow$ PT explicitly --
  every problem returned to the same idea.     \newline
  Preview Topic 4 LEHS Q\#3 and Q\#5 (asymptotes from rational
  expressions -- they come from these same domain restrictions).
  &
  Listen to the synthesis.     \newline
  One pair re-states the throughline in their own words (cold-call).     \newline
  Use the word bank if helpful: \emph{domain restriction, simplified,
  equivalent, closed, undefined}.
  &
  How will students express what they have learned?     \newline
  How does today's lesson connect to upcoming Topic 4 assessment? \\ \hline

  % ---------- EXIT TICKET ----------
  Exit Ticket (5 min)
  &
  Direct students to Page 4 (\#14 Use Appropriate Tools).
  Set timer 5 minutes.
  Collect at the door as students leave.     \newline
  Tonight: scan exits for misconceptions about
  $\frac{-6x^2+21x}{3x}$ vs.\ $-2x + 7$ at $x = 0$ -- the should-be-a-hole
  vs.\ defined-point distinction. If many students miss it, open Monday's
  Do Now with a bridge re-prompt.
  &
  Read \#14 prompt.     \newline
  Use the sentence frames if needed.     \newline
  Complete reflection prompts (``learned'' / ``liked'').     \newline
  Hand to teacher at the door.
  &
  How will students show what they've learned based on the objective?     \newline
  Are the sentence frames being used? \\ \hline
\end{xltabular}

\end{document}
